% ECONOMICS_PROBLEM: {"id": "D71-648", "title": "The midpoint satisfaction threshold in budget division", "jel_code": "D71", "source_id": "OP-0225"}
% !TeX program = xelatex
\documentclass[11pt,a4paper]{article}
\usepackage[margin=23mm]{geometry}
\usepackage{fontspec}
\setmainfont{Latin Modern Roman}
\usepackage{amsmath,amssymb,mathtools}
\usepackage{xcolor,enumitem,booktabs,longtable,array,xurl,hyperref}
\definecolor{muted}{HTML}{53636C}
\hypersetup{colorlinks=true,urlcolor=blue,linkcolor=blue}
\setlength{\parindent}{0pt}
\setlength{\parskip}{5pt}
\newcommand{\R}{\mathbb R}
\newcommand{\E}{\mathbb E}
\newcommand{\N}{\mathbb N}
\newcommand{\ind}{\mathbf 1}
\newcommand{\Var}{\operatorname{Var}}
\newcommand{\Cov}{\operatorname{Cov}}
\newcommand{\supp}{\operatorname{supp}}
\DeclareMathOperator*{\argmin}{arg\,min}
\DeclareMathOperator*{\argmax}{arg\,max}
\newcommand{\entrymeta}[3]{{\sffamily\small\color{muted}\textbf{\texttt{#1}}\quad|\quad #2\par #3\par}}
\newcommand{\Problemref}[1]{\texttt{#1}}
\newcommand{\JELref}[2]{\texttt{#2} (JEL \texttt{#1})}
\begin{document}
\section*{The midpoint satisfaction threshold in budget division}
\textbf{Problem ID:} \texttt{D71-648}\quad \textbf{JEL:} \texttt{D71}\par
% BEGIN ECONOMICS STATEMENT
\label{problem:OP-0225}
% GAME_THEORY_PROBLEM: {"local_id": "GT-095", "original_number": 95, "kind": "P", "entry_kind": "statement_only", "published": false, "review_required": true, "category": "game_theory"}
\label{game:GT-095}
{\sffamily\small\color{muted}
\noindent\textbf{ID:} GT-095. \textbf{Category:} Public allocation and computational complexity.
\textbf{Status:} P; midpoint complexity question in the inspected source version.
\par}
\subsection*{Definitions and mathematical statement}
An instance has $n\geq1$ agents, an even number $m\geq2$ of projects, and rational demands
$d_{ij}\in[0,1]$ written in binary, with $\sum_{j=1}^m d_{ij}\leq1$ for each agent $i$. A budget allocation is a vector $x\in\mathbb R_{\geq0}^{m}$ satisfying $\sum_jx_j\leq1$. Agent $i$ is satisfied on project $j$ precisely when $x_j\geq d_{ij}$. Define the decision language
\[
 \begin{split}
 \operatorname{ALL\text{-}AGENTS\text{-}SAT}_{1/2}
 =\{(d_{ij}):\ &\exists x\in\mathbb R_{\geq0}^m,
 \ \sum_jx_j\leq1,\\[-2pt]
 &\forall i\in\{1,\ldots,n\},\
       |\{j:x_j\geq d_{ij}\}|\geq m/2\}.
 \end{split}
\]
Determine its computational complexity when both $n$ and $m$ are part of the input: in particular, does it have a polynomial-time algorithm, or is it NP-complete under polynomial-time many-one reductions?

Membership in NP has a finite rational-certificate interpretation: from any successful allocation, lower each $x_j$ to the largest demand on project $j$ that it meets, or to zero if it meets none. Satisfaction is preserved, total spending does not increase, and every coordinate belongs to $\{0,d_{1j},\ldots,d_{nj}\}$.

\subsection*{Short English statement}
How hard is it to fund at least half the projects sufficiently for every agent simultaneously?

\subsection*{Variants and scope boundaries}
For odd $m$, the literal inequality with $m/2$ is equivalent to a threshold $\lceil m/2\rceil$; a floor threshold is a different problem. The even-$m$ statement avoids that ambiguity. Demands equal to zero count as satisfied even on unfunded projects. Requiring exactly exhausted budget is equivalent here by monotonicity, but requiring every demand vector to sum exactly to one restricts the input class. A complete complexity classification may be more subtle than the displayed polynomial-time/NP-complete alternatives.

\subsection*{Source relationship and status}
The model and threshold problem follow Section 1.1 and Definition 2; Section 6 identifies the $m/2$ case. The catalog links v2; v1 was the retrieved primary text used for the definition and question, so no claim about changes in v2 is made.

\subsection*{References}
\begingroup\footnotesize\raggedright
\noindent [S1] Laurent Gourv\`es, Michael Lampis, Nikolaos Melissinos, and Aris Pagourtzis, \emph{Satisfactory Budget Division} (2025), arXiv:2502.00484v1. Inspected locators: Section 1.1; Observation 1; Section 4, Definition 2; Section 6, midpoint question.
\url{https://arxiv.org/html/2502.00484v1}
\par\endgroup

\subsection*{Common conventions: Source mathematical conventions}

\textbf{Finite games.} For a finite set $A$, $\Delta(A)$ is its probability
simplex. Players choose independent mixed strategies in Nash equilibrium;
correlation is allowed only when explicitly part of the solution concept.
In a game with payoff functions $u_i$ and strategy profile $x$, an additive
$\varepsilon$ Nash equilibrium satisfies
\[
 u_i(x)\geq u_i(x_i',x_{-i})-\varepsilon
 \quad\text{for every player $i$ and every admissible deviation $x_i'$}.
\]
For numerical approximation, payoffs are normalized as locally specified.
An $\varepsilon$ payoff residual is distinct from distance at most
$\varepsilon$ to the set of exact equilibria. Point convergence, convergence
of empirical play, and convergence in distance to an equilibrium set are
also distinct.

\textbf{Computational representations.} Unless stated otherwise, a complexity
question takes rational data with binary encoding; polynomial time is measured
in total bit length. A fixed-accuracy polynomial algorithm is not necessarily
polynomial in $1/\varepsilon$, and neither is necessarily polynomial in
$\log(1/\varepsilon)$. Succinct games and value, demand, or sampling oracles
must declare their interfaces. Exact real solutions require an explicit
representation; an unspecified list of arbitrary real numbers is not a
finite computational output. A claimed algorithmic barrier is meaningful
only for its specified model and reductions.

\textbf{Dynamic games.} Histories, information, observation, and allowable
randomization are part of the model. A stationary strategy depends only on
the current state; a positional strategy in a deterministic graph game
chooses one action at each controlled vertex. Nash equilibrium at an initial
state and equilibrium after every history are different requirements.
An expectation of a pathwise $\liminf$ payoff, a $\liminf$ of expected averages,
a limiting expected average, and a uniform finite-horizon equilibrium are
different criteria. None is silently substituted for another.

\textbf{Allocation and mechanisms.} A complete allocation assigns every item
exactly once unless otherwise stated. Zero-valued goods, negative values,
capacity constraints, feasibility, lotteries, and copies of goods affect
fairness definitions and are specified locally. Dominant-strategy incentive
compatibility, Bayesian incentive compatibility, universal truthfulness,
and truthfulness in expectation impose different inequalities. Whenever
an objective ratio has zero denominator, its convention is stated or those
instances are excluded.

\textbf{Infinite-dimensional models.} Expectations, integrals, stochastic
equations, and optimization objectives require their local regularity,
integrability, filtrations, and admissible domains. A backward--forward
mean-field system is a boundary-value problem; it is not an initial-value
semiflow without an additional construction. Long-time, large-population,
and vanishing-noise limits require an order of limits and a convergence
topology. If a minimum need not be attained, the objective is its infimum
or supremum and the approximation requirement is stated separately.
% END ECONOMICS STATEMENT
\end{document}
